\chapter{数值解法——多步法}

\intro{
	我们只能看到自己已经理解的东西。
	
	\rightline{——Johann Wolfgang von Goethe}
}

\section{多步法的基本思想}

与单步法仅使用 $t_n$ 处的信息不同，\textbf{线性多步法}利用之前若干节点的信息来构造 $y_{n+1}$，其一般形式为：
\begin{myformula}
	\sum_{j=0}^{k} \alpha_j y_{n+j} = h \sum_{j=0}^{k} \beta_j f_{n+j}
\end{myformula}
其中 $f_{n+j} = f(t_{n+j}, y_{n+j})$，$k$ 为步数，$\alpha_k \neq 0$，通常取 $\alpha_k = 1$。

\begin{itemize}
	\item 若 $\beta_k = 0$，格式为\textbf{显式}（Adams-Bashforth 型）；
	\item 若 $\beta_k \neq 0$，格式为\textbf{隐式}（Adams-Moulton 型）。
\end{itemize}

\section{基于数值积分的构造方法}

将微分方程 $y' = f(t, y)$ 在区间 $[t_n, t_{n+1}]$ 上积分：
\[
y(t_{n+1}) - y(t_n) = \int_{t_n}^{t_{n+1}} f(t, y(t)) \, dt
\]

用插值多项式 $P(t)$ 近似 $f(t, y(t))$，则：
\[
y_{n+1} = y_n + \int_{t_n}^{t_{n+1}} P(t) \, dt
\]

\subsection{Adams-Bashforth 显式公式}

使用 $t_{n-k+1}, \ldots, t_n$ 这 $k$ 个节点上的已知值构造插值多项式。以下列出前四种公式：

\begin{myTheorem}{Adams-Bashforth 公式（AB-k）}
	\begin{enumerate}
		\item AB-1（即 Euler 法）：$y_{n+1} = y_n + h f_n$
		\item AB-2：$y_{n+1} = y_n + \dfrac{h}{2} (3 f_n - f_{n-1})$
		\item AB-3：$y_{n+1} = y_n + \dfrac{h}{12} (23 f_n - 16 f_{n-1} + 5 f_{n-2})$
		\item AB-4：$y_{n+1} = y_n + \dfrac{h}{24} (55 f_n - 59 f_{n-1} + 37 f_{n-2} - 9 f_{n-3})$
	\end{enumerate}
\end{myTheorem}

\begin{proof}[AB-2 的推导]
	使用 $t_{n-1}, t_n$ 两节点构造线性插值：
	\[
	P(t) = f_n + \frac{f_n - f_{n-1}}{h} (t - t_n)
	\]
	对 $P(t)$ 从 $t_n$ 到 $t_{n+1}$ 积分：
	\begin{align*}
		\int_{t_n}^{t_{n+1}} P(t) \, dt &= \int_{0}^{h} \left[ f_n + \frac{f_n - f_{n-1}}{h} \tau \right] d\tau \\
		&= h f_n + \frac{f_n - f_{n-1}}{h} \cdot \frac{h^2}{2} = \frac{h}{2} (3 f_n - f_{n-1})
	\end{align*}
	故 $y_{n+1} = y_n + \frac{h}{2}(3f_n - f_{n-1})$。
\end{proof}

\begin{remark}
	$k$ 步 AB 方法具有 $k$ 阶精度。但随着步数增加，稳定性区间会缩小。
\end{remark}

\subsection{Adams-Moulton 隐式公式}

同样用 $k$ 个节点构造插值多项式，但将 $t_{n+1}$ 也作为插值节点。

\begin{myTheorem}{Adams-Moulton 公式（AM-k）}
	\begin{enumerate}
		\item AM-1（即隐式 Euler）：$y_{n+1} = y_n + h f_{n+1}$
		\item AM-2（即梯形法）：$y_{n+1} = y_n + \dfrac{h}{2} (f_{n+1} + f_n)$
		\item AM-3：$y_{n+1} = y_n + \dfrac{h}{12} (5 f_{n+1} + 8 f_n - f_{n-1})$
		\item AM-4：$y_{n+1} = y_n + \dfrac{h}{24} (9 f_{n+1} + 19 f_n - 5 f_{n-1} + f_{n-2})$
	\end{enumerate}
\end{myTheorem}

\begin{remark}
	同阶的 AM 方法比 AB 方法精度更高（$k$ 步 AM 方法为 $k+1$ 阶），且具有更好的稳定性。但由于是隐式，需要结合显式预测使用。
\end{remark}

\section{预测-校正方法（Adams-Bashforth-Moulton）}

最流行的做法是用 AB 方法作预测、AM 方法作校正，形成\textbf{PECE} 格式（Predict-Evaluate-Correct-Evaluate）：

\begin{myformula}
	\begin{cases}
		\text{P:} \quad y_{n+1}^{(0)} = y_n + \dfrac{h}{24} (55 f_n - 59 f_{n-1} + 37 f_{n-2} - 9 f_{n-3}) \\[8pt]
		\text{E:} \quad f_{n+1}^{(0)} = f(t_{n+1}, y_{n+1}^{(0)}) \\[8pt]
		\text{C:} \quad y_{n+1} = y_n + \dfrac{h}{24} (9 f_{n+1}^{(0)} + 19 f_n - 5 f_{n-1} + f_{n-2}) \\[8pt]
		\text{E:} \quad f_{n+1} = f(t_{n+1}, y_{n+1})
	\end{cases}
\end{myformula}

\textbf{PECE 模式}的步骤：
\begin{enumerate}
	\item \textbf{预测 (Predict)}：用 AB-4 显式公式计算 $y_{n+1}^{(0)}$；
	\item \textbf{评估 (Evaluate)}：计算 $f_{n+1}^{(0)}$；
	\item \textbf{校正 (Correct)}：用 AM-4 隐式公式修正 $y_{n+1}$；
	\item \textbf{评估 (Evaluate)}：计算新的 $f_{n+1}$ 供下一步使用。
\end{enumerate}

\section{多步法的启动问题}

多步法不能自启动——需要在开始时用同阶的单步法（通常为 RK4）计算出足够多的起始值。启动过程如下：

\begin{enumerate}
	\item 由初始值 $y_0$ 出发，用 RK4 步进 $k-1$ 步，生成 $y_1, \ldots, y_{k-1}$；
	\item 同时计算对应的 $f_1, \ldots, f_{k-1}$；
	\item 此后用 $k$ 步线性多步法继续推进。
\end{enumerate}

\section{线性多步法的阶条件}

\begin{theorem}{多步法的阶条件}
	线性多步法 $\sum_{j=0}^{k} \alpha_j y_{n+j} = h \sum_{j=0}^{k} \beta_j f_{n+j}$ 具有 $p$ 阶精度的充要条件是系数满足：
	\[
	\begin{aligned}
		\sum_{j=0}^{k} \alpha_j &= 0 \\
		\sum_{j=0}^{k} \left( \frac{j^q}{q!} \alpha_j - \frac{j^{q-1}}{(q-1)!} \beta_j \right) &= 0, \quad q = 1, 2, \ldots, p
	\end{aligned}
	\]
\end{theorem}

\begin{proof}
	将精确解 $y(t_{n+j})$ 在 $t_n$ 处 Taylor 展开：
	\[
	y(t_{n+j}) = \sum_{m=0}^{\infty} \frac{(jh)^m}{m!} y^{(m)}(t_n)
	\]
	而 $f(t_{n+j}, y(t_{n+j})) = y'(t_{n+j}) = \sum_{m=0}^{\infty} \frac{(jh)^m}{m!} y^{(m+1)}(t_n)$。
	
	代入多步法格式，合并同次幂的 $h^m$ 项系数。为使局部截断误差为 $O(h^{p+1})$，需 $h^0, h^1, \ldots, h^p$ 各项系数为零，即得上列阶条件。
\end{proof}

\section{向后微分公式（BDF）}

对于\textbf{刚性}（stiff）问题，需要具有良好稳定性性质的方法。\textbf{向后微分公式}（Backward Differentiation Formula, BDF）直接对 $y(t)$ 进行插值，然后令插值多项式的导数在 $t_{n+k}$ 处等于 $f_{n+k}$。

$k$ 步 BDF 格式为：
\begin{myformula}
	\sum_{j=0}^{k} \alpha_j y_{n+j} = h \beta_k f_{n+k}
\end{myformula}

\begin{myTheorem}{常用 BDF 公式}
	\begin{enumerate}
		\item BDF-1（隐式 Euler）：$y_{n+1} - y_n = h f_{n+1}$
		\item BDF-2：$y_{n+2} - \dfrac{4}{3} y_{n+1} + \dfrac{1}{3} y_n = \dfrac{2}{3} h f_{n+2}$
		\item BDF-3：$y_{n+3} - \dfrac{18}{11} y_{n+2} + \dfrac{9}{11} y_{n+1} - \dfrac{2}{11} y_n = \dfrac{6}{11} h f_{n+3}$
	\end{enumerate}
\end{myTheorem}

\begin{remark}
	BDF 方法对刚性方程有极佳的稳定性——BDF-1 和 BDF-2 是 $A$-稳定的，BDF-3 到 BDF-6 是 $A(\alpha)$-稳定的。MATLAB 的 \texttt{ode15s} 求解器即基于 BDF 方法，适用于求解刚性 ODE。
\end{remark}

\section{常用 MATLAB 求解器总结}

\begin{table}[H]
	\centering
	\caption{MATLAB ODE 求解器一览}
	\begin{tabular}{clcc}
		\toprule
		求解器 & 方法 & 类型 & 适用场景 \\
		\midrule
		\texttt{ode45} & Dormand-Prince (4,5) 对 & 显式 & 非刚性，首选 \\
		\texttt{ode23} & Bogacki-Shampine (2,3) 对 & 显式 & 轻度刚性，低精度 \\
		\texttt{ode113} & Adams-Bashforth-Moulton PECE & 多步 & 非刚性，高精度 \\
		\texttt{ode15s} & BDF (1--5) & 隐式 & 刚性问题 \\
		\texttt{ode23s} & 修正 Rosenbrock & 单步 & 刚性，低精度 \\
		\texttt{ode23t} & 梯形法 & 隐式 & 中度刚性 \\
		\texttt{ode23tb} & TR-BDF2 & 隐式 & 刚性问题 \\
		\bottomrule
	\end{tabular}
\end{table}

\section{数值方法的实现框架}

以一阶方程组 $\dot{\mathbf{y}} = \mathbf{f}(t, \mathbf{y})$ 为例，数值求解的通用流程为：

\begin{enumerate}
	\item 定义右端函数 $\mathbf{f}(t, \mathbf{y})$（匿名函数或函数文件）；
	\item 设定时间区间 $[t_0, T]$ 与初值 $\mathbf{y}_0$；
	\item 调用求解器获得解向量 $[\mathbf{t}, \mathbf{Y}]$；
	\item 后处理：绘图分析、参数优化。
\end{enumerate}

\begin{example}{MATLAB 求解指数衰减方程}
	以 $\dot{y} = -a y, \; y(0) = y_0$ 为例：
	\begin{verbatim}
	a = 0.5;  tspan = [0 20];  y0 = 1;
	odefun = @(t,y) -a * y;
	[t, y] = ode45(odefun, tspan, y0);
	plot(t, y, 'b-', 'LineWidth', 2);
	\end{verbatim}
\end{example}

\begin{example}{Python 显式 Euler 法的实现}
	求解 $\dot{y} = -a y$，展示不同步长的影响：
	\begin{verbatim}
	def euler_solve(a, dt, t_max, y0):
	    t, y = 0.0, y0
	    t_list, y_list = [t], [y]
	    while t < t_max:
	        t += dt
	        y += dt * (-a * y)
	        t_list.append(t)
	        y_list.append(y)
	    return t_list, y_list
	\end{verbatim}
\end{example}
